\section{Evidence for the Broader Behavioral Findings}
\label{app:behavior-evidence}

The following analyses and experiments belong to the present investigation.
They retain different sampling frames, deployment aliases, and decision rules.
Table~\ref{tab:behavior-map} gives the full evidence map, including the
persistence results developed after Section~\ref{sec:differences}.
Table~\ref{tab:evidence-anchors} supplies quantitative anchors for the broader
behavioral findings. Its accompanying inventory maps each quantitative row to
a saved aggregate, field, and source hash. These tables report existing
results; they do not introduce a new analysis or a common validation rule.

\begin{table}[t]
\centering\small
\begin{tabular}{p{0.22\linewidth}p{0.34\linewidth}p{0.35\linewidth}}
\toprule
Behavior & Supported finding & Boundary of evidence\\
\midrule
Answer revelation; reasoning elicitation & Defaults predict new problems and tested new expressions & Binary events; probabilities and ranks need not be invariant\\
Assistance directness & Multi-method, cross-task behavioral tendency & Complete semantic-screen decision; fixed panel\\
Feeling/difficulty acknowledgement & Default and state increments on new problems with original wording & State increment does not transfer to new expressions\\
Response complexity & Cross-task semantic signature & Rated steps/information, not measured learner load\\
Expressed confidence & Cross-context consistency in reported confidence & Exploratory archive evidence; not correctness or calibration\\
Affiliation inventory & All five models choose the affiliative option on the default items & Shared choices; open-response differences and limits reported separately\\
Safety-boundary facets & Several within-task recurring patterns & Tasks do not support one permissiveness axis\\
Other candidate dimensions & Mixed or insufficient full-chain evidence & Failed checks are retained, not evidence of universal absence\\
\bottomrule
\end{tabular}
\caption{Detailed evidence map across the investigation.
Behavioral screens and prospective tests address different questions; inventory
choices describe responses to fixed options. Quantitative anchors and measurement distinctions are reported in this appendix.}
\label{tab:behavior-map}
\end{table}

\begin{table}[t]
\centering\small
\begin{tabular}{llrr}
\toprule
ID & Behavior & Statistic A & Statistic B\\
\midrule
E1 & Assistance directness & 0.629 & 0.647\\
E2 & Response complexity & 0.663 & 0.677\\
E3 & Semantic elicitation & 0.475 & 0.468\\
E4 & Expressed confidence & 0.852 & 0.886\\
E5 & Revision propensity & 0.537 & 0.472\\
E6 & Unanswerable abstention & 0.346 & 0.705\\
E7 & Attack success & 0.806 & 0.943\\
E8 & SATA incorrect inclusion & 0.761 & 0.714\\
E9 & Next-step actionability & -0.111 & -0.771\\
E10 & Open affiliation & 0.900 & 0.718\\
\bottomrule
\end{tabular}
\caption{Quantitative anchors for the supplementary behavioral findings. For E1--E9, A is cross-context ICC(3,1) and B is median pairwise model-rank Spearman correlation. For E10, A is education/non-education rank correlation and B is default/irrelevant-change rank correlation. These measure different forms of recurrence; they are not comparable trait-strength scores or new significance tests.}
\label{tab:evidence-anchors}
\end{table}


\subsection{Semantic Screening and Candidate Dimensions}

The semantic panel rates visible responses from six historical configurations
on eight dimensions using three model coders. Assistance directness spans
hint/question-only assistance to a complete solution; cognitive load spans a
manageable next step to many simultaneous steps or concepts. Thus the latter
is an output rubric, not a learner measurement. Ratings are examined for coder
agreement, variation, cross-task consistency, and links to observable actions
and prompt changes. Coders' response-level agreement and models' cross-task
consistency are separate quantities. Directness alone satisfies the complete
saved behavioral-disposition rule. Cognitive load satisfies the narrower
semantic-signature criteria; action/quality evidence does not support its
promotion to a fully validated disposition. No learning-efficacy claim follows
from the directness decision either.

The decision rule was fixed after partial coder results had been inspected,
before the remaining coder reveal; it is not a fully outcome-blind
preregistration. It constrains subsequent decisions without converting the
whole semantic screen into independent prospective confirmation.

\paragraph{Complete screening rule.}
The hierarchy requires all preceding tiers; passing one component alone does
not validate a dimension. Reliable measurement requires coder ICC(3,$k$)
$\geq .60$, consensus-score SD $\geq .50$, neither endpoint above 80\%,
minimum leave-one-coder model-profile Spearman $\geq .70$, and no
leave-one-coder prompt-effect sign reversal among headline model--task
contrasts. Cross-task signature status additionally requires a positive lower
95\% bootstrap bound for model variance in every sampled benchmark and either
cross-task ICC(3,1) $\geq .50$ or median pairwise-task Spearman $\geq .50$.
Whole-panel held-out-task attribution must exceed $1/6$ chance; this shared
check is not a dimension-specific validation test.

The strongest tier further requires at least one of three corroborations:
(1) adding the full semantic panel to a transparent quality predictor yields
positive mean leave-one-model-out AUC gain with no negative model fold, and
the dimension separates independently inferred human actions in its anchored
direction; (2) paired prompt movement agrees with the independently trained
human-action result and keeps its direction under every single-coder exclusion;
or (3) on 40 LongTutor histories \citep{li2026longtutor}, the dimension predicts
human-gold diagnosis beyond exact history, with within-outcome BH-adjusted
$q<.05$ under the specified history permutation test. Only directness passes
the complete hierarchy, through criterion (2). Semantic elicitation passes that
corroboration but fails the preceding cross-task tier.

The thresholds were fixed on 19 August 2026 after 637 of 1,080 planned coder
annotations, before the third coder and LongTutor results were available.
An implementation clarification at 640 annotations applied the variance check
to every benchmark; the BH amendment at 725 annotations preceded all LongTutor
coding. These timings are disclosed because this was a staged decision rule,
not an outcome-blind preregistration. Self-family residual differences or slate-position
ranges above .25 scale points flag interpretation concerns; they do not replace
the tier decisions.

\begin{table}[t]
\centering\small
\begin{tabular}{lccc}
\toprule
Semantic dimension & Reliable & Cross-task & Complete\\
\midrule
Assistance directness & Pass & Pass & Pass\\
Elicitation & Pass & Fail & Fail\\
Autonomy support & Fail & Fail & Fail\\
Affective warmth & Fail & Fail & Fail\\
Diagnostic specificity & Fail & Fail & Fail\\
Personalization & Fail & Fail & Fail\\
Response complexity & Pass & Pass & Fail\\
Epistemic caution & Fail & Fail & Fail\\
\bottomrule
\end{tabular}
\caption{Cumulative decisions for all eight semantic dimensions. Failure
at an earlier tier prevents promotion, irrespective of later component tests.
The binary prospective elicitation event is a separate measure.}
\label{tab:screen-decisions}
\end{table}

The retained eight-dimensional record includes elicitation, autonomy support,
affective warmth, diagnostic specificity, personalization, and epistemic
caution, with their failed components. A separate four-coder frozen-split
panel examines instructional agency, relational communion, and next-step
actionability in matched standard/hard and Socratic settings. Agency's failed
pilot gate is not reversed by stronger formal agreement. Communion has no
model variation in the constrained Socratic setting and overlaps with warmth.
Actionability fails cross-task stability (E9), despite an instruction-induced
increase for every tested model in both matched task variants. Outcome-aware
source-grouped reanalysis is reported in Appendix~\ref{app:archive-reanalysis};
it does not retroactively change these decisions. In particular, the semantic
elicitation score (E3) and the subsequent binary reasoning-invitation endpoint
are not identical measures.

\subsection{Confidence, Revision, and Safety Facets}

These are outcome-aware exploratory analyses of already scored archive
outputs. The common six historical configurations are MiniMax-M3,
MiniMax-M2.7, GLM-5.2, DeepSeek-V4-Pro, Doubao-Seed-2.0-Pro, and Qwen3.5-4B.
The epistemic analysis uses recorded confidence, revision, correctness, and
abstention fields rather than newly assigning personality ratings. Model means
are compared across four source contexts for confidence and revision and five
unanswerability categories for abstention. ICC(3,1) and median pairwise rank
correlation each use a 0.50 descriptive screening threshold. Confidence passes
both (E4), revision only the ICC check (E5), and abstention only the rank check
(E6). These finite-panel screens do not identify latent factors or make reported
confidence a correctness guarantee.

Safety analysis compares attack success across five risk categories and SATA
incorrect inclusion across ten scenario-language cells (E7--E8). Their
cross-task model-rank correlation is -0.486, opposite the proposed common
permissiveness direction. Educational-refusal style also has recurring
within-task structure, but one category-pair rank correlation is only 0.086;
SATA omission has mixed consistency. These are different abilities or policy
facets, not a single validated personality axis. The adversarial scorer is
MiniMax-M3 for all candidates, including itself, and is not independently
validated across judge families.

\subsection{Targeted Affiliation Pilot}

The prospective pilot contains 280 generation calls and 144 blinded coding
calls, with twelve open conflict scenarios. Five historical configurations
participate: MiniMax-M3, MiniMax-M2.7, GLM-5.2, DeepSeek-V4-Pro, and
Doubao-Seed-2.0-Lite. Three model coders score open affiliation. The fixed gates
require coder ICC at least 0.70, cross-domain and irrelevant-change rank
correlations at least 0.70, absolute mean irrelevant displacement at most
0.25, and default between-model standard deviation at least 0.15. All pass:
coder ICC is 0.931, the two rank correlations are in E10, displacement is
-0.050, and model standard deviation is 0.445. With only five configurations,
the two exact two-sided permutation p-values are 0.091 and 0.174. This is a
local gate-based result, not population-level significance or broad invariance.

High-minus-low affiliation instructions shift the mean by 1.689 on the 1--5
scale, with a positive direction in all five models. Default rank correlation
is 0.900 under high affiliation and 0.200 under low affiliation. Self-report
correlation with open behavior is -0.200; default forced-choice responses have
no model variance. Open affiliation correlates 0.900 with surface warmth.
Thus locally recurring open behavior is supported, while general-trait
convergence and separation from warm expression are not established.

\subsection{Semantic Prompt Effects, Ordering, and Attribution}
\label{app:prompt-response}

This secondary analysis was designed after inspecting the results. It uses paired generic and
pedagogy responses already rated for six historical models: MiniMax-M3,
MiniMax-M2.7, GLM-5.2, DeepSeek-V4-Pro, Doubao-Seed-2.0-Pro, and Qwen3.5-4B.
There are 79 standard and 40 hard MathDial contexts, each with both prompt
conditions and all six models. The eight semantic dimensions yield 11,424
response--dimension rows, not 11,424 independent replies. Analyses use the
saved three-annotator consensus. They do not revise the recorded decisions about the dimensions
or constitute independent prospective testing.

\paragraph{Prompt shift relative to default spread.}
For each task and dimension, the absolute mean paired pedagogy-minus-generic
shift is divided by the full range of model means under the generic prompt.
Table~\ref{tab:prompt-scale} reports the three dimensions that passed the
measurement reliability criterion, with 5,000 context-bootstrap intervals. These are
ratios of descriptive point estimates, not tests that prompts universally
matter more than model choice. A ratio near one means that the mean shift
under this instruction is comparable to the observed baseline range.

\begin{table}[t]
\centering\small
\begin{tabular}{llr}
\toprule
Task & Dimension & Shift / default range [95\% CI]\\
\midrule
Standard & Response complexity & 1.111 [0.785, 1.401]\\
Standard & Semantic elicitation & 1.148 [0.946, 1.437]\\
Standard & Assistance directness & 1.000 [0.830, 1.204]\\
Hard & Response complexity & 0.848 [0.585, 1.214]\\
Hard & Semantic elicitation & 0.987 [0.777, 1.301]\\
Hard & Assistance directness & 1.011 [0.758, 1.361]\\
\bottomrule
\end{tabular}
\caption{Prompt shifts relative to default model spread for the rated behavioral dimensions. CI denotes confidence interval. Intervals are descriptive context-bootstrap intervals, not simultaneous tests across dimensions.}
\label{tab:prompt-scale}
\end{table}

\paragraph{Response heterogeneity and scale boundaries.}
Let $s$ and $p$ denote the generic and pedagogy scores. For a pooled increase,
the adjusted change is $(p-s)/(5-s)$; for a pooled decrease, it is
$(s-p)/(s-1)$. Only contexts
where every model has nonzero room to change in the tested direction are retained. The retained
standard/hard counts are 41/19 for elicitation, 58/31 for directness, and 36/16
for complexity. Model labels are permuted within paired contexts 5,000 times;
Benjamini--Hochberg correction covers eight dimensions within each task.
Adjusted $q$ values for elicitation are 0.0003/0.0016 and for directness
0.0003/0.0155 (standard/hard); complexity is 0.0003/0.7023. Thus only
elicitation and directness among the reliably measured dimensions retain
heterogeneity in both variants after adjustment. The exclusion of boundary
contexts changes the analysis subset, and normalization addresses one
floor/ceiling explanation rather than every possible source of heterogeneity.

\paragraph{Changes in model ordering.}
Table~\ref{tab:prompt-order-example} illustrates the Doubao/GLM directness
comparison. Across the six models, directness ordering reverses for 5/15
comparable pairs in standard and 6/14 in hard contexts; pairs tied in either
condition are excluded. These counts describe sample-mean crossings, not separate
significant pairwise reversals. The generic-to-pedagogy Spearman correlations
are 0.371 and 0.203; their exact two-sided permutation $p$ values, 0.497 and
0.711, do not establish either preserved ordering or its absence.

\begin{table}[t]
\centering\small
\begin{tabular}{llrrr}
\toprule
Task & Model & Generic & Pedagogy & Paired change [95\% CI]\\
\midrule
Standard & Doubao-Seed-2.0-Pro & 3.90 & 1.58 & -2.32 [-2.62, -2.01]\\
Standard & GLM-5.2 & 2.97 & 1.87 & -1.10 [-1.38, -0.84]\\
Hard & Doubao-Seed-2.0-Pro & 3.98 & 2.10 & -1.88 [-2.30, -1.43]\\
Hard & GLM-5.2 & 3.08 & 2.48 & -0.60 [-1.00, -0.17]\\
\bottomrule
\end{tabular}
\caption{An illustrative change in directness ordering. Higher scores indicate more direct assistance. Intervals describe each model's paired prompt shift, not the difference between models or a test of rank reversal.}
\label{tab:prompt-order-example}
\end{table}

\paragraph{Residual model attribution across instructions.}
A standardized multinomial logistic classifier predicts the six model
identities from eight item-centered semantic scores, training under one
instruction and testing under the other. All eight dimensions are included;
classification success does not validate them individually as traits.
Pooled generic-to-pedagogy accuracy is 0.301 [0.264, 0.335]; the reverse is
0.297 [0.266, 0.326], compared with chance 0.167. The four cross-task,
cross-prompt directions range from 0.263 to 0.300, with lower interval bounds
above chance. Intervals use 2,000 test-context bootstrap samples conditional
on the fitted classifier. Pooled transfers reuse paired contexts across conditions;
cross-task transfers compare the two benchmark variants rather than a new
prospectively sampled source population. This is supplementary evidence of
joint behavioral discrimination, distinct from the action-probability
predictions fixed before response collection in Section~\ref{sec:stability}. It does not establish immutable
model identity or an internal personality mechanism.


\subsection{Prompt Control and the Action-Selection Boundary}

Factorial action instructions are additionally tested under randomized order
and two new clause paraphrases. Question-first and correct-answer-revelation
controls pass the two-wording rules; a literal encouragement marker fails one
wording's threshold and is not a measure of semantic warmth. Ordinary student
requests for the opposite action do not override the configured policy in
those tested contrasts. This supports action control, not adaptive selection.

A prospective two-stage trial uses 96 contexts, five configurations, two
selector wordings, both counterfactual executors, and a concurrent single-pass
baseline: 2,400 generation cells. Specified ASK and EXPLAIN actions are realized
at 0.998 and 1.000 under the trial's question-based binary metric. Selector
match to an existing observed-teacher binary reference is 0.519 and 0.500;
composed-minus-single-pass match is -0.058 and -0.077. Both selection wordings
fail the registered joint target. An observed teacher move is not uniquely
optimal, and question-based action realization is not a quality measure.
The trial constrains naive system-decomposition claims without becoming the
main contribution of the educational-personality investigation.
